% Generated from locale/tr/content/first-order-logic/completeness/introduction.tex; do not hand-edit.


\iftag{FOL}
      {\olfileid[tr]{fol}{com}{int}}
      {\olfileid[tr]{pl}{com}{int}}

\olsection{Giriş}

Tamlık teoremi, mantık hakkındaki en temel sonuçlardan biridir.
Birbirine denk olduklarını ispatlayacağımız iki biçimde ifade edilir.
İlk biçimi, semantik sonuç ile !!{derivation} sistemimiz arasındaki
ilişkiye dair temel bir şey söyler: !!a{sentence}~$!A$, bir
!!{sentence}s kümesi~$\Gamma$'dan çıkıyorsa, $\Gamma \Proves !A$
olduğunu ortaya koyan !!a{derivation} da vardır. Dolayısıyla
!!{derivation} sistemi, gerçekte çıkmayan şeyleri ispatlamaksızın
olabileceği kadar güçlüdür.

İkinci biçimi bir model varlığı sonucu olarak ifade edilebilir: her
tutarlı !!{sentence}s kümesi sağlanabilirdir. Tutarlılık ispat
kuramsal bir kavramdır; !!{derivation} sistemimizin belirli
!!{derivation}s üretemediğini söyler. Fakat sırf~$\Gamma$'dan belirli
türde !!{derivation}s bulunmuyor diye
\iftag{FOL}{!!a{structure}~$\Struct{M}$}{!!{valuation}{~$\pAssign{v}$}
ve $\iftag{FOL}{\Sat{M}}{\pSat{v}}{\Gamma}$ olduğunu} kim güvence altına
alabilir? Tamlık teoremi ilk kez ispatlanmadan önce---hatta bugün
kullandığımız !!{derivation} sistemleri ortaya çıkmadan önce---büyük
Alman matematikçi David Hilbert, matematiksel kuramların tutarlılığının
bu kuramların konu ettiği nesnelerin varlığını güvence altına aldığı
görüşündeydi. Gottlob Frege'ye yazdığı bir mektupta bunu şöyle dile
getirmişti:
\begin{quote}
  Keyfî olarak verilen aksiyomlar, bütün sonuçlarıyla birlikte birbiriyle
  çelişmiyorsa doğrudurlar ve aksiyomların tanımladığı şeyler vardır. Benim
  için doğruluğun ve varlığın ölçütü budur.
\end{quote}
Frege buna şiddetle karşı çıktı. Tamlık teoreminin ikinci biçimi,
aksiyomlar tutarlıysa en azından hepsini doğru kılan \emph{bir}
\iftag{FOL}{!!{structure}}{!!{valuation}} var olduğu anlamında
Hilbert'in haklı olduğunu gösterir.

Tamlık teoreminin---daha doğrusu ispatının---önemli olmasının tek
nedenleri bunlar değildir. Ayrı ayrı ele alacağımız bir dizi önemli
sonucu da vardır. Örneğin $\Gamma \Proves !A$ olduğunu gösteren her
!!{derivation} sonludur ve bu nedenle~$\Gamma$ içindeki
!!{sentence}s arasından yalnızca sonlu sayıda olanı kullanabilir.
Tamlık teoremi uyarınca bundan şu sonuç çıkar: $!A$,~$\Gamma$'nın bir
sonucuysa zaten~$\Gamma$'nın sonlu bir alt kümesinin sonucudur. Buna
\emph{tıkızlık} denir. Buna denk olarak, $\Gamma$'nın her sonlu alt
kümesi tutarlıysa $\Gamma$'nın kendisi de tutarlı olmalıdır.

Tıkızlık teoremi, !!{derivation}s üzerinden dolanan bir yolla tamlık
teoreminden çıksa da \emph{tamlık teoreminin ispatını} kullanarak onu
doğrudan ortaya koymak da mümkündür. Çünkü ispat, belirli bir özelliğe
---tutarlılığa---sahip bir !!{sentence}s kümesinden yola çıkar ve bu
kümeden belirli özellikleri olan !!a{structure} kurar (burada yapı söz
konusu kümeyi sağlar). Hemen hemen aynı kuruluş, tutarlı kümeler yerine
``sonlu olarak sağlanabilir'' !!{sentence}s kümelerinden başlanarak
tıkızlığı doğrudan ortaya koymak için kullanılabilir.
\iftag{FOL}{Bu kuruluş başka sonuçlar da verir; örneğin sağlanabilir her
  !!{sentence}s kümesinin sonlu ya da !!{denumerable}s bir modeli
  vardır. (Bu sonuca Löwenheim--Skolem teoremi denir.) Genel olarak
  !!{structure}s kurmak için !!{sentence}s kümelerinden yararlanma
  yöntemi mantıkta,
  hatta bazen felsefede bile sıkça kullanılır.}{}

